\subsection*{M222\quad 相对论变质量火箭的四动量平衡与双曲正切火箭方程}
质量随燃料喷出而减小的理想直线火箭，初始静质量$M_0$、终质量$M_f$，喷气在火箭瞬时静止系中的后向速度大小恒定$u<c$；忽略外力和辐射。
(1) 用瞬时共动系能量和动量守恒推导微分形式的速度增量。
(2) 将速度复合写成快速度，积分得到火箭末速度。
(3) 给出低速喷气与光子火箭极限。
(4) 给定目标$v_f$反求燃料质量比，并检查燃料比趋一的极限。
\par\textbf{逐问解析}
(1) 在某一瞬时火箭静止系内，发射尾流静质量$dm_e>0$、速率$u$，其能量和动量大小为$\gamma_u dm_ec^2$与$\gamma_udm_eu$，$\gamma_u=(1-u^2/c^2)^{-1/2}$。发射后火箭质量为$M+dM$（$dM<0$），微小前向速度$dv$。能量守恒到一阶给$-dM=\gamma_udm_e$，动量守恒给$M\,dv=\gamma_udm_eu$，故$dv=-u\,dM/M$；这里$dv$专指共动系的微小相对速度，而非固定实验室速度差。
(2) 用快速度$\eta=\operatorname{artanh}(v/c)$处理速度复合：局域增量$d\eta=dv/c=-(u/c)dM/M$。从$M_0$积分至$M_f$，若火箭初始静止则$\eta_f=(u/c)\ln(M_0/M_f)$，最终$v_f=c\tanh[(u/c)\ln(M_0/M_f)]$。以快速度而非线性速度相加，保证全部质量比下$v_f<c$。
(3) 若$|\eta_f|\ll1$，$\tanh\eta_f\simeq\eta_f$，回复经典齐奥尔科夫斯基公式$v_f\simeq u\ln(M_0/M_f)$。对极限$u\uparrow c$，有$v_f/c=\tanh\ln(M_0/M_f)=(M_0^2-M_f^2)/(M_0^2+M_f^2)$，即理想光子火箭方程；有限质量比始终不能达到光速。
(4) 定义目标快速度$\eta_f=\operatorname{artanh}(v_f/c)$，反求$M_0/M_f=\exp[(c/u)\operatorname{artanh}(v_f/c)]$。若质量比趋一，指数趋零，所以$v_f\to0$；若$u$固定且$v_f\to c^-$，质量比发散。该模型将尾流能量包含在火箭质量损失内，因此不能再无端把排气静质量误认为$-dM$而遗漏$\gamma_u$。
\par\medskip
\subsection*{M223\quad 平面限制性三体问题的三角Lagrange点与Routh稳定阈值}
规范化圆形限制性三体问题中，主星质量$1-\mu$和伴星质量$\mu$（$0<\mu\le1/2$）在旋转系固定为$(-\mu,0)$、$(1-\mu,0)$，旋转角速度为1；无质量粒子受引力及惯性力作用。
(1) 建立有效势$\Omega(x,y)$及运动方程。
(2) 验证三角平衡点$L_{4,5}$坐标。
(3) 推导三角点二维线性扰动特征多项式。
(4) 求谱稳定的临界质量比并解释等号共振极限。
\par\textbf{逐问解析}
(1) 定义$r_1=\sqrt{(x+\mu)^2+y^2}$、$r_2=\sqrt{(x-1+\mu)^2+y^2}$，有效势$\Omega(x,y)=(x^2+y^2)/2+(1-\mu)/r_1+\mu/r_2$。旋转系方程为$\ddot x-2\dot y=\Omega_x$、$\ddot y+2\dot x=\Omega_y$，其中二阶速度项是Coriolis加速度，势中的二次项包含离心作用。
(2) 在等边三角形构型，两主星到无质量粒子的距离均为1，取$x_*=1/2-\mu$、$y_*=\pm\sqrt3/2$。于是$\Omega_x=x_*-(1-\mu)(x_*+\mu)-\mu(x_*-1+\mu)=0$，$\Omega_y=y_*-(1-\mu)y_*-\mu y_*=0$，因此两点确为平衡位置。
(3) 对平衡点取$\xi,\eta\propto e^{st}$，势函数Hessian满足$\Omega_{xx}=3/4$、$\Omega_{yy}=9/4$、$\Omega_{xy}=\pm(3\sqrt3/4)(1-2\mu)$。线性矩阵行列式为$(s^2-\Omega_{xx})(s^2-\Omega_{yy})+(2s-\Omega_{xy})(2s+\Omega_{xy})=0$，故$s^4+[4-\Omega_{xx}-\Omega_{yy}]s^2+\Omega_{xx}\Omega_{yy}-\Omega_{xy}^2=0$，化为$s^4+s^2+27\mu(1-\mu)/4=0$。
(4) 令$z=s^2$，两根为$z_\pm=[-1\pm\sqrt{1-27\mu(1-\mu)}]/2$。二者均严格负实的条件是$27\mu(1-\mu)<1$，所以$0<\mu<\mu_R=[1-\sqrt{23/27}]/2\simeq0.0385209$谱稳定；超过阈值根成复共轭，出现增长模。在$\mu=\mu_R$处有纯虚重根，必须单独判定Jordan共振和非线性稳定性，不能简单归为严格稳定。
\par\medskip
